\documentclass{article}

\usepackage{amsmath}
\usepackage{amssymb}
\usepackage[utf8]{inputenc}
\usepackage[T1]{fontenc}
\usepackage[polish]{babel}

\title{Wielomiany}
\date{2026}

\begin{document}
\maketitle

\section{Definicja i stopień}
Niech $K$ będzie ciałem. Wielomianem jednej zmiennej $x$ o współczynnikach
w~$K$ nazywamy wyrażenie postaci
\[
    P(x)=a_0+a_1x+a_2x^2+\cdots+a_nx^n,
    \qquad a_0,\ldots,a_n\in K.
\]
Wielomian zerowy oznaczamy przez $0$ i przyjmujemy, że jego stopień nie jest
określony (czasami zapisuje się $\deg 0=-\infty$). Jeżeli $P\neq 0$ oraz
$a_n\neq 0$, to liczbę $n$ nazywamy stopniem wielomianu $P$ i oznaczamy
przez $\deg P$. Współczynnik $a_n$ nazywa się współczynnikiem wiodącym,
a $a_0$ wyrazem wolnym.

Wielomian można traktować zarówno jako formalne wyrażenie, jak i funkcję
$K\to K$ daną wzorem $x\mapsto P(x)$. Dla ciał nieskończonych, na przykład
$\mathbb{R}$ lub $\mathbb{C}$, dwa wielomiany są równe wtedy i tylko wtedy,
gdy mają te same współczynniki. Nad ciałem skończonym różne wielomiany mogą
wyznaczać tę samą funkcję.

Zbiór wszystkich wielomianów o współczynnikach w $K$ oznaczamy przez
$K[x]$. Zbiór wielomianów stopnia co najwyżej $n$ oznaczamy przez
\[
    K_{[n]}[x]=\{0\}\cup\{P\in K[x]\colon 1\leq\deg P\leq n\}.
\]
Dla $K=\mathbb{R}$ często piszemy krótko $W_{[n]}$.

\section{Działania na wielomianach}
Jeżeli
\[
    P(x)=\sum_{k=0}^{m}a_kx^k,\qquad
    Q(x)=\sum_{k=0}^{n}b_kx^k,
\]
to dodawanie i mnożenie przez skalar definiujemy współczynnikowo:
\[
    (P+Q)(x)=\sum_{k\geq 0}(a_k+b_k)x^k,\qquad
    (\alpha P)(x)=\sum_{k=0}^{m}(\alpha a_k)x^k,
\]
przy czym brakujące współczynniki traktujemy jako równe zero. Iloczyn
wielomianów ma postać
\[
    (PQ)(x)=\sum_{k=0}^{m+n}c_kx^k,
    \qquad c_k=\sum_{j=0}^{k}a_jb_{k-j}.
\]
W szczególności, dla niezerowych wielomianów:
\[
    \deg(P+Q)\leq\max\{\deg P,\deg Q\},\qquad
    \deg(PQ)=\deg P+\deg Q.
\]
Pierwsza nierówność może być ostra, gdy wyrazy najwyższego stopnia się
skasują.

\section{Dzielenie z resztą}
Jeżeli $A,B\in K[x]$ oraz $B\neq 0$, to istnieją jednoznacznie wyznaczone
wielomiany $Q,R\in K[x]$ takie, że
\[
    A=BQ+R,\qquad R=0\ \text{lub}\ \deg R<\deg B.
\]
Wielomian $Q$ nazywamy ilorazem, a $R$ resztą z dzielenia $A$ przez $B$.
Jeżeli $R=0$, mówimy, że $B$ dzieli $A$ i zapisujemy $B\mid A$.

Szczególnym przypadkiem jest dzielenie przez $x-a$. Wtedy
\[
    P(x)=(x-a)Q(x)+P(a).
\]
Jest to twierdzenie o reszcie z dzielenia.

\section{Pierwiastki i czynniki liniowe}
Element $a\in K$ nazywamy pierwiastkiem wielomianu $P$ wtedy i tylko wtedy,
gdy $P(a)=0$. Z twierdzenia o reszcie otrzymujemy twierdzenie Bézouta:
\[
    P(a)=0\quad\Longleftrightarrow\quad (x-a)\mid P(x).
\]
Jeżeli $a$ jest pierwiastkiem wielomianu $P$, to możemy więc zapisać
\[
    P(x)=(x-a)Q(x).
\]
Jeżeli
\[
    P(x)=(x-a)^rQ(x),\qquad Q(a)\neq 0,
\]
to mówimy, że $a$ jest pierwiastkiem krotności $r$. Pierwiastek krotności
co najmniej $2$ nazywamy pierwiastkiem wielokrotnym.

Jeżeli $P$ ma stopień $n$ i rozkłada się nad $K$ na czynniki liniowe, to
\[
    P(x)=a_n(x-a_1)(x-a_2)\cdots(x-a_n),
\]
gdzie pierwiastki są uwzględnione wraz z krotnościami. W szczególności
niezerowy wielomian stopnia $n$ ma co najwyżej $n$ różnych pierwiastków.

\subsection{Sprzężone pierwiastki}
Jeżeli $P\in\mathbb{R}[x]$ oraz $z\in\mathbb{C}$ jest jego pierwiastkiem,
to liczba sprzężona $\overline{z}$ również jest pierwiastkiem. Wynika to z
równości
\[
    P(\overline{z})=\overline{P(z)}.
\]
Zatem
\[
    P(z)=0\quad\Longrightarrow\quad P(\overline{z})=0.
\]
Pierwiastki zespolone niebędące rzeczywistymi występują więc w parach
$z,\overline{z}$. Co więcej, oba pierwiastki mają taką samą krotność.

\section{Twierdzenie zasadnicze algebry}
Każdy niezerowy wielomian stopnia $n\geq 1$ o współczynnikach zespolonych ma
co najmniej jeden pierwiastek zespolony. W konsekwencji każdy wielomian
$P\in\mathbb{C}[x]$ stopnia $n$ rozkłada się jednoznacznie z dokładnością do
kolejności czynników w postaci
\[
    P(x)=a_n\prod_{j=1}^{n}(x-a_j).
\]
Dla wielomianów rzeczywistych rozkład na czynniki liniowe może nie istnieć
w $\mathbb{R}[x]$, ale każdy taki wielomian rozkłada się na iloczyn
czynników liniowych i nierozkładalnych czynników kwadratowych.

\section{Wielomiany jako przestrzeń liniowa}
Zbiór $K_{[n]}[x]$ z naturalnym dodawaniem wielomianów i mnożeniem przez
skalary jest przestrzenią liniową nad $K$. Jego bazą kanoniczną jest
\[
    \mathcal{B}=(1,x,x^2,\ldots,x^n),
\]
a wymiar wynosi
\[
    \dim K_{[n]}[x]=n+1.
\]
Współrzędne wielomianu
$P(x)=a_0+a_1x+\cdots+a_nx^n$ w tej bazie są równe
\[
    [P]_{\mathcal{B}}=
    \begin{bmatrix}a_0\\a_1\\\vdots\\a_n\end{bmatrix}.
\]
Przestrzeń $K[x]$ wszystkich wielomianów ma bazę
$(1,x,x^2,\ldots)$ i jest nieskończenie wymiarowa.

\end{document}
